\section{Вибір криптографічного механізму}

\subsection{Ситуація 1. Захист електронного документа}

\textbf{Умова:} Необхідно забезпечити можливість перевірити, що електронний документ був створений певним власником закритого ключа та не був модифікований після підписання.

\textbf{Запропонований криптографічний механізм:}

\textbf{Можливий алгоритм або схема:}

\textbf{Обґрунтування вибору:}

\subsection{Ситуація 2. Формування спільного секретного ключа}

\textbf{Умова:} Alice та Bob повинні сформувати спільний секретний ключ, обмінюючись інформацією через недовірену мережу.

\textbf{Запропонований криптографічний механізм:}

\textbf{Можливий алгоритм або схема:}

\textbf{Обґрунтування вибору:}

\subsection{Ситуація 3. Конфіденційна передача великого файла}

\textbf{Умова:} Необхідно передати файл розміром 2 ГБ у конфіденційному вигляді. Відкритий ключ отримувача відомий відправнику.

\textbf{Запропонований криптографічний механізм:}

\textbf{Можливий алгоритм або схема:}

\textbf{Обґрунтування вибору та переваги гібридної схеми:}

\subsection{Ситуація 4. Безпечне зберігання паролів}

\textbf{Умова:} Сервер повинен безпечно зберігати паролі користувачів та перевіряти їх під час автентифікації.

\textbf{Чи потрібна асиметрична криптографія:}

\textbf{Запропонований механізм:}

\textbf{Обґрунтування вибору:}

\subsection{Ситуація 5. Автентичність оновлень програмного забезпечення}

\textbf{Умова:} Виробник випускає оновлення програмного забезпечення для великої кількості пристроїв. Створювати справжні оновлення може лише виробник, а кожен пристрій повинен мати можливість перевірити їх автентичність.

\textbf{Запропонований криптографічний механізм:}

\textbf{Можливий алгоритм або схема:}

\textbf{Розподіл ключів між виробником та пристроями:}

\textbf{Обґрунтування вибору:}